%%FILE blueprint/src/macros/common.tex
% In this file you should put all LaTeX macros and settings to be used both by
% the pdf version and the web version.
% This should be most of your macros.

% The theorem-like environments defined below are those that appear by default
% in the dependency graph. See the README of leanblueprint if you need help to
% customize this.
% The configuration below use the theorem counter for all those environments
% (this is what the [theorem] arguments mean) and never resets it.
% If you want for instance to number them within chapters then you can add
% [chapter] at the end of the next line.
\newtheorem{theorem}{Theorem}[section]
\newtheorem{proposition}[theorem]{Proposition}
\newtheorem{lemma}[theorem]{Lemma}
\newtheorem{corollary}[theorem]{Corollary}

\theoremstyle{definition}
\newtheorem{definition}[theorem]{Definition}

\renewcommand{\d}{\, \mathrm{d}}

\newcommand{\R}{\mathbb{R}}
\newcommand{\C}{\mathbb{C}}
\newcommand{\N}{\mathbb{N}}
\newcommand{\Z}{\mathbb{Z}}
\newcommand{\B}{\mathcal{B}}
\newcommand{\tor}{\mathbb{T}}

\newcommand{\incset}{\mathrm{incidenceSet}}
\newcommand{\dartset}{\mathrm{dartSet}}
\newcommand{\loopset}{\mathrm{loopSet}}
\newcommand{\incother}{\mathrm{inc}\_\mathrm{other}}
\newcommand{\const}{\mathbf{1}}
\newcommand{\isloop}{\mathrm{IsLoop}}
\newcommand{\islink}{\mathrm{IsLink}}
\newcommand{\cc}{\mathrm{ConnComp}}
\newcommand{\ccm}{\mathrm{cc}}
\newcommand{\symm}{\mathrm{rev}}
\newcommand{\spt}{\mathrm{start}}
\newcommand{\ept}{\mathrm{end}}
\newcommand{\edge}{\mathrm{edge}}
\newcommand{\perm}[1]{\mathfrak{S}_{#1}}
\newcommand{\de}{\mathfrak{d}}
\newcommand{\length}{\mathrm{length}}
\newcommand{\diam}{\mathrm{diam}}
\newcommand{\ediam}{\mathrm{ediam}}
\newcommand{\dist}{\mathrm{dist}}
\newcommand{\edist}{\mathrm{edist}}
\newcommand{\Prob}{\mathrm{Prob}_{N,r}}
\newcommand{\ProbG}{\mathrm{Prob}_{N,\de}^{G}}
\newcommand{\expec}{\mathbb{E}_{N,r}}
\newcommand{\tree}{\mathbb{T}_\de}
\newcommand{\free}{\mathbb{F}_r}
\newcommand{\word}{\mathbf{W}_r}
\newcommand{\matrices}{\mathcal{M}_N(\C)}
\newcommand{\trN}[1][N]{\mathrm{tr}_{#1}}
\newcommand{\TrN}[1][N]{\mathrm{Tr}_{#1}}
\newcommand{\supp}{\mathrm{supp} \,}
\newcommand{\testfunc}[1][\R, \C]{\mathcal{D}(#1)}
\newcommand{\distrib}[1][\R, \C]{\mathcal{D}'(#1)}
\newcommand{\smooth}[1][\R, \C]{\mathcal{E}(#1)}
\newcommand{\csdistrib}[1][\R, \C]{\mathcal{E}'(#1)}
\newcommand{\RS}{\hat{\mathbb{C}}}
\newcommand{\nonco}{\mathcal{M}_D(\C) \otimes \C \langle x_1, \ldots, x_{r} \rangle}

%%FILE blueprint/src/macros/web.tex
% In this file you should put macros to be used only by
% the web version. Of course they should have a corresponding
% version in macros/print.tex.
% Typically the printed version could have more fancy decorations.
% This will probably be a very short file.
%%FILE blueprint/src/macros/print.tex
% In this file you should put macros to be used only by
% the printed version. Of course they should have a corresponding
% version in macros/web.tex.
% Typically the printed version could have more fancy decorations.
% This should be a very short file.
%
% This file starts with dummy macros that ensure the pdf
% compiler will ignore macros provided by plasTeX that make
% sense only for the web version, such as dependency graph
% macros.


% Dummy macros that make sense only for web version.
\newcommand{\lean}[1]{}
\newcommand{\discussion}[1]{}
\newcommand{\leanok}{}
\newcommand{\mathlibok}{}
\newcommand{\notready}{}
% Make sure that arguments of \uses and \proves are real labels, by using invisible refs:
% latex prints a warning if the label is not defined, but nothing is shown in the pdf file.
% It uses LaTeX3 programming, this is why we use the expl3 package.
\ExplSyntaxOn
\NewDocumentCommand{\uses}{m}
 {\clist_map_inline:nn{#1}{\vphantom{\ref{##1}}}%
  \ignorespaces}
\NewDocumentCommand{\proves}{m}
 {\clist_map_inline:nn{#1}{\vphantom{\ref{##1}}}%
  \ignorespaces}
\ExplSyntaxOff